\chapter{Programmcode}
\label{anhang_programmcode}
\lstinputlisting[language=Matlab, caption={Skript für die Autocodegenerierung und Nachbearbeitung}, label=lst_build_skript_nachbearbeitung_ganz]{build_release.m} 
\chapter{Abkürzungen Matlab-Modell}
%https://www.tablesgenerator.com/

% Please add the following required packages to your document preamble:
% \usepackage[table,xcdraw]{xcolor}
% Beamer presentation requires \usepackage{colortbl} instead of \usepackage[table,xcdraw]{xcolor}
% \usepackage{longtable}
% Note: It may be necessary to compile the document several times to get a multi-page table to line up properly
\begin{longtable}{|l|l|l|l|}
\hline
\rowcolor[HTML]{C0C0C0} 
\textbf{Abkürzung} & \textbf{Bezeichnung}         & \textbf{Zusatzinfo}    & \textbf{neu?} \\ \hline
\endfirsthead
%
\endhead
%
Akku               & Akkumulator                  &                         & Ja            \\ \hline
Aktv               & aktivieren, aktiv            &                         & Nein          \\ \hline
Anm                & Anlaufmodus                  &                         & Nein          \\ \hline
Aus                & aus, ausschalten             &                         & Nein          \\ \hline
Blk                & blockieren, Blockade         &                         & Ja            \\ \hline
Btr                & Betrieb                      &                         & Nein          \\ \hline
bus                & Bus-Signal                   & Spezieller Variablentyp & Nein          \\ \hline
Dek                & dekrementieren, Dekrement    &                         & Nein          \\ \hline
dm                 & Drehmoment                   & physikalische Größe     & Nein          \\ \hline
Ein                & ein, einschalten             &                         & Nein          \\ \hline
Elk                & Elektrik, Elektronik         &                         & Ja            \\ \hline
Erk                & Erkennung                    &                         & Nein          \\ \hline
flg                & Flag                         & Spezieller Variablentyp & Nein          \\ \hline
Ink                & inkrementieren, Inkrement    &                         & Nein          \\ \hline
Lim                & Limit                        &                         & Ja            \\ \hline
lst                & Leistung                     & physikalische Größe     & Ja            \\ \hline
Max                & Maximal                      &                         & Nein          \\ \hline
Mot                & Motor, Maschine              &                         & Nein          \\ \hline
Mtst               & Motorsteuerung               &                         & Nein          \\ \hline
n                  & Drehzahl                     & physikalische Größe     & Nein          \\ \hline
Nsr                & not Safety relevant          &                         & Ja            \\ \hline
num                & Nummer, Zahl, Ziffer, Anzahl & Spezieller Variablentyp & Nein          \\ \hline
Ok                 & In Ordnung, Ok               &                         & Ja            \\ \hline
Red                & reduzieren, Reduzierung      &                         & Ja            \\ \hline
spng               & Spannung                     & physikalische Größe     & Nein          \\ \hline
Sr                 & Safety relevant              &                         & Ja            \\ \hline
stm                & Strom                        & physikalische Größe     & Nein          \\ \hline
Stst               & Stillstand                   &                         & Ja            \\ \hline
t                  & Zeit                         & physikalische Größe     & Nein          \\ \hline
tmp                & Temperatur                   & physikalische Größe     & Nein          \\ \hline
Zael               & Zähler                       &                         & Nein          \\ \hline
Zn                 & zu niedrig                   &                         & Ja            \\ \hline
\caption{Abkürzungsverzeichnis Matlab - Modell}
\label{tab_Namen_Matlab}\\
\end{longtable}